\documentclass[11pt,a4paper]{article}

\usepackage[utf8]{inputenc}
\usepackage[T1]{fontenc}
\usepackage[spanish,es-tabla]{babel}
\usepackage{amsmath,amssymb,amsthm,mathtools}
\usepackage{booktabs}
\usepackage{enumitem}
\usepackage{geometry}
\usepackage{setspace}
\usepackage{hyperref}
\usepackage{microtype}

\geometry{margin=1in}
\onehalfspacing

\hypersetup{
  pdftitle={Nota formal: Invariantes y robustez de Systemic Tau},
  pdfauthor={Johel Padilla-Villanueva}
}

\theoremstyle{definition}
\newtheorem{definition}{Definici\'on}[section]
\newtheorem{remark}{Observaci\'on}[section]
\theoremstyle{plain}
\newtheorem{proposition}{Proposici\'on}[section]
\newtheorem{theorem}{Teorema}[section]
\newtheorem{corollary}{Corolario}[section]
\newtheorem{conjecture}{Conjetura}[section]
\theoremstyle{remark}
\newtheorem{scholium}{Scholium}[section]

\newcommand{\Res}{\mathrm{Res}}
\newcommand{\RECD}{\mathrm{RECD}}
\newcommand{\TC}{\mathrm{TC}}
\newcommand{\Syn}{\mathrm{Syn}}
\newcommand{\Surp}{\mathrm{Surp}}
\newcommand{\excess}{\mathrm{excess3}}
\newcommand{\Rpair}{\mathcal{R}_{\mathrm{pair}}}
\newcommand{\KL}{D_{\mathrm{KL}}}
\newcommand{\ts}{\tau_s}

\title{\textbf{Nota formal}\\[0.4em]
\large Invariantes, covarianza y robustez\\
de Systemic Tau y de las medidas de conjunci\'on\\[0.6em]
\normalsize (Punto 5 del roadmap; asume Notas 1--4)}

\author{
Johel Padilla-Villanueva\\
\textit{Department of Environmental Health, University of Puerto Rico}\\
\texttt{johelpadilla@gmail.com}
}

\date{2026-08-01 \quad$\cdot$\quad v0.1}

\begin{document}
\maketitle

\begin{abstract}
Se clasifican las transformaciones bajo las cuales $\ts$, $\Phi_\ell$,
$\Res_{\mathrm{pair}}$, $\excess$ y $T_{\RECD}$ son invariantes,
covariantes de forma controlada, o inestables. Se proponen versiones
can\'onicas / normalizadas y se discuten l\'imites asint\'oticos en $d$
(n\'umero de variables) y $T$ (longitud / ventana). El objetivo es
cerrar el punto~5 del roadmap: robustez te\'orica de la medida, no solo
comportamiento en un ensamble sint\'etico fijo.
\end{abstract}

\tableofcontents
\vspace{1em}

\section{Objetos}
\label{sec:obj}

Trabajamos con: proceso $X_t\in\mathbb{R}^d$; s\'imbolos Bandt--Pompe
$S_t\in\Sigma^d$, $\Sigma=S_m$;
\begin{align}
\ts(t)&=\frac{2}{d(d-1)}\sum_{i<j}\tau_K\bigl(r^{(i)}_t,r^{(j)}_t\bigr),\\
\Phi_1,\Phi_2,&\quad
\Res_{\mathrm{pair}},\ \excess,\quad
\Delta\RECD=\sum_\ell\alpha_\ell(\lambda)\widehat{M}_\ell.
\end{align}

\section{Invariancias exactas}
\label{sec:exact}

\begin{theorem}[Invariancia mon\'otona componente a componente]
\label{thm:mono}
Si $Z_t^{(i)}=f_i\bigl(X_t^{(i)}\bigr)$ con cada $f_i$
\emph{estrictamente creciente}, entonces:
\begin{enumerate}[label=(\roman*)]
  \item los s\'imbolos $S_t$ de $Z$ coinciden con los de $X$;
  \item $\Phi_1,\Phi_2,\TC,\Res_{\mathrm{ind}},\Res_{\mathrm{pair}},
  \excess$ (sobre $S$) son id\'enticos;
  \item si $\ts$ se calcula sobre rangos/series ordinales equivalentes,
  $\ts$ es id\'entico (Kendall es invariante mon\'otono por coordenada
  cuando se aplica la misma $f_i$ a cada serie en la ventana).
\end{enumerate}
\end{theorem}

\begin{proof}
Bandt--Pompe depende solo del orden relativo dentro de cada retardo;
Kendall $\tau_K$ depende solo de concordancias de rangos.
\end{proof}

\begin{theorem}[Permutaci\'on de canales]
\label{thm:perm}
Si $\sigma\in S_d$ reordena coordenadas,
$\ts$, $\Phi_\ell$, $\Res_{\mathrm{pair}}$, $\excess$ son invariantes.
$T_{\RECD}$ tambi\'en, si $\lambda$ se construye de forma sim\'etrica en
los canales.
\end{theorem}

\begin{theorem}[Traslaci\'on temporal]
\label{thm:shift}
Si $Y_t=X_{t+t_0}$, todos los funcionales de ventana se trasladan:
$G_Y(t)=G_X(t+t_0)$. No hay preferencia de origen temporal.
\end{theorem}

\begin{proposition}[Invariancia de escala global por canal]
\label{prop:scale}
$X^{(i)}\mapsto a_i X^{(i)}+b_i$ con $a_i>0$ cae en
Thm.~\ref{thm:mono}. Si alg\'un $a_i<0$ (reflexi\'on), los s\'imbolos
Bandt--Pompe se conjugan por la inversi\'on de orden: $\Phi_1$ (igualdad
de s\'imbolos) se preserva solo si \emph{todas} las coordenadas se
reflejan o ninguna en un sentido compatible; en general
\textbf{no} hay invariancia bajo reflexiones mixtas.
\end{proposition}

\section{Covariancias controladas}
\label{sec:cov}

\begin{proposition}[Cambio de embedding $(m,\tau)$]
\label{prop:embed}
$\Phi_\ell$ y $\Res$ viven en $\Sigma=S_m$ distinto al cambiar $m$. No
son num\'ericamente comparables entre $m$ distintos sin renormalizaci\'on.
Protocolo: fijar $(m,\tau)$ a priori (defaults $m=3$, $\tau=1$) o
reportar familia indexada por $m$.
\end{proposition}

\begin{proposition}[Longitud de ventana $w$]
\label{prop:w}
Estimadores plug-in de entrop\'ia / $\Res$ / $\excess$ son sesgados en
$w$ peque\~no y variables en $w$ grande (no estacionariedad). Covarianza:
sesgo $\downarrow$ y varianza de estimaci\'on $\downarrow$ al crecer $w$
bajo estacionariedad local; resoluci\'on temporal $\downarrow$.
\end{proposition}

\begin{proposition}[Dependencia en $d$]
\label{prop:d}
\begin{enumerate}[label=(\roman*)]
  \item $\ts$ promedia $O(d^2)$ pares: concentraci\'on mejora con $d$
  si los pares son d\'ebilmente dependientes; un par outlier se diluye.
  \item $\Res_{\mathrm{pair}}$ en $\Sigma^d$ sufre la maldici\'on de
  dimensionalidad: $|\Sigma|^d=(m!)^d$ crece r\'apido; para $m=3$, $d=4$
  es factible; $d\gg 1$ exige proyecciones a subconjuntos o estimadores
  regularizados.
  \item $\Phi_1,\Phi_2$ normalizados por $\binom{d}{2}$ permanecen en
  $[0,1]$ y son comparables entre $d$ si la geometr\'ia de pares es
  homog\'enea.
\end{enumerate}
\end{proposition}

\begin{proposition}[Motor $\alpha(\lambda)$]
\label{prop:alpha_cov}
$T_{\RECD}$ \textbf{no} es invariante bajo cambio de familia $\alpha$
admisible (Nota~3). Solo lo son los $\widehat{M}_\ell$ crudos. Toda
comparaci\'on de relojes entre estudios exige declarar $\alpha$ y
$\lambda$.
\end{proposition}

\section{Versiones can\'onicas / normalizadas}
\label{sec:canonical}

\begin{definition}[Lectura can\'onica de abundancia L3]
\label{def:can_A3}
Reportar el triple
\begin{equation}
\bigl(\Syn,\ \Surp,\ \excess\bigr)
\quad\text{y, si es factible,}\quad
\Res_{\mathrm{pair}}=\KL(P\|P^{(2)}),
\end{equation}
sin reponderar. Esta es la forma can\'onica de \emph{abundancia}.
\end{definition}

\begin{definition}[Share can\'onica]
\label{def:can_f}
$f_\ell$ solo con motor admisible preespecificado y $\lambda$ en la
jerarqu\'ia de fallback (Nota~3). Etiqueta:
$f_\ell^{(\alpha,\lambda)}$.
\end{definition}

\begin{definition}[Normalizaci\'on por nulo]
\label{def:nullnorm}
\begin{equation}
\Delta A_3
:= A_3^{\mathrm{obs}}-A_3^{\mathrm{null}},
\end{equation}
donde el nulo es (i)~independencia (shuffle de canales) o
(ii)~surrogados que preservan bipolares. $\Delta A_3$ es la versi\'on
can\'onica robusta a offset de estimaci\'on.
\end{definition}

\begin{definition}[$\ts$ en bandas]
\label{def:bands}
Comparar regímenes por ocupaci\'on de bandas
($|\ts|<\tau_{\mathrm{ch}}$, $\ts\ge 0.50$, etc.) adem\'as del valor
medio: m\'as estable a monotom\'ias afines del cue.
\end{definition}

\section{Asint\'otica}
\label{sec:asymp}

\begin{conjecture}[Ley de grandes n\'umeros local]
\label{conj:lln}
Bajo mezcla suficientemente r\'apida del proceso $S_t$ y $w\to\infty$
con $w=o(T)$, los estimadores plug-in de $\TC$, $\excess$ y
$\Res_{\mathrm{pair}}$ convergen en probabilidad a sus valores
poblacionales en la medida estacionaria (o en la media erg\'odica
local).
\end{conjecture}

\begin{conjecture}[Concentraci\'on de $\ts$]
\label{conj:conc_tau}
Si los $\tau_K^{(ij)}$ son d\'ebilmente dependientes,
$\mathrm{Var}(\ts)=O(1/d^2)+O(\rho/d)$ con $\rho$ correlaci\'on media
entre pares; $\ts$ se concentra al crecer $d$.
\end{conjecture}

\begin{proposition}[L\'imite $d=2$]
\label{prop:d2}
Para $d=2$: $\Syn\equiv 0$, L3 fuerte vac\'io, $\ts=\tau_K$ del \'unico
par, $\Phi_1\in\{0,1\}$ por tiempo. El paradigma de profundidad exige
$d\ge 3$.
\end{proposition}

\section{Robustez pr\'actica (lista de control)}
\label{sec:checklist}

\begin{enumerate}[label=\textbf{R\arabic*.}]
  \item Fijar $(m,\tau,w,p_{\min},f_{\min})$ a priori.
  \item Verificar invariancia mon\'otona en un smoke test
  ($X\mapsto\log(1+X)$, etc.).
  \item No comparar $\excess$ entre $m$ distintos sin nota.
  \item Para $d$ grande: L3 en subconjuntos / sketch, no joint completo.
  \item Separar $A_3$ de $f_3^{(\alpha,\lambda)}$.
  \item Publicar nulo bipolar adem\'as de nulo de independencia.
  \item Si $\lambda$ falla el design goal: $\alpha\equiv 1$.
\end{enumerate}

\section{Canon a\~nadido (C29--C34)}
\label{sec:canon}

\begin{enumerate}[label=\textbf{C\arabic*.}]
  \setcounter{enumi}{28}
  \item Invariancia mon\'otona estricta por canal (Thm.~\ref{thm:mono}).
  \item Invariancia bajo permutaci\'on de canales.
  \item No invariancia gen\'erica bajo reflexiones mixtas ni bajo cambio
  de $m$.
  \item Formas can\'onicas: triple Syn/Surp/excess3; $\Delta A_3$ vs nulo;
  $f^{(\alpha,\lambda)}$.
  \item $d\ge 3$ para L3 fuerte; $d=2$ degenera Syn.
  \item $T_{\RECD}$ depende del motor: no comparar relojes sin declarar
  $\alpha,\lambda$.
\end{enumerate}

\vspace{1.5em}
\noindent\rule{\textwidth}{0.4pt}\\[0.5em]
\begin{center}
\em Fin Nota formal punto 5 v0.1.
\end{center}

\begin{thebibliography}{99}
\bibitem{Padilla2026note1}
J.~Padilla-Villanueva. Notas formales 1--4. Fundaci\'on RECD, 2026.
\bibitem{BandtPompe2002}
C.~Bandt and B.~Pompe. Permutation entropy.
\emph{Phys.\ Rev.\ Lett.} 88, 174102 (2002).
\end{thebibliography}

\end{document}
